% Chapter/section outline for the YANI paper.
% Bullets are the intended content of each section, not final prose.
% The abstract lives beside this file in paper/abstract.md.

\section{Introduction}
\begin{itemize}
  \item Activation, not transport, is what sets maintenance access, cooling
        and waste classification in a fusion plant: the inventory problem is
        the one the design decisions actually hang on.
  \item The problem statement is flux-field general -- fusion, fission,
        accelerator and spallation sources, radioisotope sources -- and the
        paper is written that way, with fusion as the driving case.
  \item Landscape of existing inventory codes (FISPACT-II, ORIGEN, ALARA,
        Serpent and OpenMC depletion), what each assumes about group
        structure, network truncation and coupling to transport.
  \item Gaps this work targets: multigroup pre-collapse, fixed-depth network
        truncation, ground-state-only branching, installation and licensing
        friction, and post-processing burden.
  \item Contributions of the paper, and a roadmap of the sections.
\end{itemize}

\section{Theory}
\begin{itemize}
  \item Bateman equations with reaction terms; production and loss by both
        decay and neutron reaction; definition of $N_i$, $\lambda$,
        $\sigma\phi$.
  \item Assembly of the burnup matrix $A$: which coefficients come from the
        transmutation network and which from the spectrum collapse.
  \item One step as a matrix exponential $\mathbf{N}(t)=e^{At}\mathbf{N}(0)$;
        why the exponential is the cheap part and the coefficients are not.
  \item CRAM at order 48 (Pusa 2010, 2015): rational approximation, pole
        structure, accuracy against analytic chains, and the stiffness
        argument for choosing it over explicit integration.
  \item Sparse LU with a symbolic factorization reused across poles; sparsity
        pattern of a real network; cost model per step.
  \item Beginning-of-step reaction rates: what a single stepper does and does
        not resolve, and the error incurred over a long high-flux step.
  \item Continuous-energy collapse of $\sigma\phi$ against a user spectrum,
        contrasted with a library pre-collapsed on someone else's weighting
        spectrum; resonance self-shielding caveats.
  \item Energy-dependent isomeric branching folded against the spectrum
        rather than applied as an evaluation-average ratio.
\end{itemize}

\section{Network construction and reduction}
\begin{itemize}
  \item What a transmutation network is here: decay constants, decay modes
        and branching, reaction products by MT, fission yields, isomeric
        overlays.
  \item Reachability from the starting nuclides: why a fixed reaction/decay
        depth both drops real nuclides and loads unused ones.
  \item The provable upper bound on the density each candidate nuclide can
        reach given the irradiation time and rates, and the floor it is
        tested against.
  \item Effect on network size, data volume loaded and solve time for
        representative fusion materials (steel, tungsten, lithium ceramics).
  \item Verification that the reduction is conservative: comparison against
        an unreduced solve.
\end{itemize}

\section{Nuclear data}
\begin{itemize}
  \item Pipeline: evaluated ENDF-6 in, NJOY processing, a Rust ENDF parser,
        Apache Arrow IPC out; one conversion, many runs.
  \item The on-disk representation: per-nuclide section objects, typed
        columns, LZ4 compression, memory-mapped reads, and why this replaces
        parsing fixed-width text at startup.
  \item Schema: cross-section sections, transmutation networks, decay data,
        fission yields, branching overlays, and how they are versioned.
  \item Library coverage and what each library can supply
        (ENDF/B-VIII.1, JEFF-4.0, TENDL-2025, TENDL-2017, FENDL-3.2d);
        which have decay data and fission yields of their own.
  \item Independently sourced subsections: mixing reactions from one library
        with decay data and yields from another as a first-class
        configuration, and the provenance record that goes with it.
  \item Processing-code sensitivity: the IAEA-NDS NJOY fork versus upstream
        for FENDL, quantified (URR probability tables, MT 301 heating,
        MT 444 damage), and why the build refuses the wrong flavour.
  \item Distribution: publication per library, on-demand fetch of only the
        sections a calculation needs, disk cache, and the in-memory cache
        that is topped up rather than rebuilt.
  \item Integrity gate: readback validation, compression checks, and
        rebuild-to-rebuild comparison that separates a re-encode from a
        change in the numbers.
\end{itemize}

\section{Software design and implementation}
\begin{itemize}
  \item A single Rust workspace shared with the YAMC Monte Carlo transport
        code; crate layout (solver, decay, conversion, bindings, WASM) and
        what is shared rather than duplicated.
  \item Python bindings: thin wrappers over Rust, typed stubs as the single
        source of truth for the API and the generated documentation.
  \item Caching and memory: process-wide caches keyed on data source, top-up
        on a widened request, and the behaviour of a script that loops over
        many materials or schedule variants.
  \item Parallelism and where the time goes: per-material independence, and
        the cost split between network build, rate collapse and the
        exponential.
  \item Packaging: wheels for Linux, macOS and Windows on x86 and ARM, with
        no compiler required at install; the metadata-package arrangement
        that keeps one source of truth for the API.
  \item Licensing: a fully permissive stack, and what that allows downstream.
  \item Testing and CI: unit tests, Python-level tests, doctested
        documentation, and regression gates on numerical results.
\end{itemize}

\section{User interface}
\begin{itemize}
  \item Materials: composition by element or nuclide, natural abundance
        expansion, density and volume, and the bundled PNNL Compendium
        (PNNL-15870 Rev.~2) material library.
  \item Spectra: arbitrary group structures, named structures as a
        convenience, normalisation, and passing a transport tally straight
        in.
  \item Schedules: pulses and cooldowns, per-pulse spectra so a campaign
        whose spectrum changes shape (DD then DT) is expressed directly,
        and step-size guidance.
  \item Nuclear data settings: the five module-level settings, per-nuclide
        overrides, local files and downloaded keywords in one calculation.
  \item Results: inventories, activity, decay heat, contact dose and photon
        lines, each available totalled or by nuclide, and the plotting
        helpers that keep post-processing out of user scripts.
  \item Worked example end to end, with the code that produces each figure.
\end{itemize}

\section{Derived quantities}
\begin{itemize}
  \item Activity and decay heat from the inventory; treatment of stable
        nuclides and of deep traces.
  \item Decay photon line spectra as discrete energy/intensity pairs rather
        than a binned response, and why lines matter for nuclide
        identification and for feeding transport.
  \item Contact dose: the half-space slab estimate, self-attenuation, the
        build-up factor, and why no distance or volume enters.
  \item Dose quantities: absorbed dose in air and ICRP-116 effective dose;
        the NIST $\mu/\rho$ and $\mu_{en}/\rho$ tabulations used, and
        agreement with the FISPACT-II methodology.
  \item Known omission: no bremsstrahlung from decay electrons, and the
        materials where that reads low.
  \item Coupling to transport: emitting the photon source in the form a
        coupled photon transport run consumes directly, and the shutdown
        dose rate workflow that follows from it.
\end{itemize}

\section{Verification and validation}
\begin{itemize}
  \item Analytic verification: single decay chains, secular and transient
        equilibrium, and constant-rate irradiation with closed-form
        solutions.
  \item Solver verification: CRAM order and step-size convergence, sparse
        versus dense factorization, and mass conservation.
  \item Code-to-code comparison against established inventory codes on
        common materials and spectra, with the data library held fixed so
        the comparison is of methods rather than of evaluations.
  \item Validation against open activation benchmarks, CoNDERC among them:
        activity, decay heat and gamma emission after irradiation, with
        C/E ratios per nuclide and per cooling time.
  \item Contact dose cross-check against an independent implementation of
        the same methodology.
  \item Sensitivity of the results to the choice of library and of NJOY
        flavour, separating method error from evaluation spread.
  \item Reproducibility: the V\&V suite as runnable scripts pinned to
        library versions, run in CI.
\end{itemize}

\section{Performance}
\begin{itemize}
  \item Timing breakdown for a representative case: data load, network
        build, rate collapse, factorization, exponential.
  \item Scaling with network size, number of schedule steps and number of
        materials; cold cache versus warm cache.
  \item Startup cost of the columnar data format against a text-parsing
        baseline, and the effect of fetching only the sections needed.
  \item Memory footprint, and the effect of the network reduction on it.
  \item Native versus WebAssembly throughput on the same case.
  \item Cross-platform and cross-architecture results (Windows, macOS,
        Linux; x86 and ARM).
\end{itemize}

\section{Browser deployment}
\begin{itemize}
  \item The same Rust solver compiled to WebAssembly: one codebase, not a
        reduced browser edition.
  \item Architecture of the page: self-contained, solve running
        client-side, nuclear data as the only fetch, nothing uploaded.
  \item Equivalence with the installed wheel: identical continuous-energy
        data and identical answers, demonstrated on the benchmark cases.
  \item What this changes for teaching, review and reproducibility: a
        calculation that a reader can rerun with no install and no account.
  \item Export of a browser session as a Python script, as the path from
        exploration to a scripted workflow.
\end{itemize}

\section{Applications}
\begin{itemize}
  \item Foil activation as a minimal case, including a metastable product
        that dominates the activity long after its ground state has decayed.
  \item First-wall and blanket structural materials under a DT spectrum:
        inventory, decay heat and contact dose against cooling time.
  \item A phased campaign with a changing spectrum, showing what a single
        averaged spectrum would have missed.
  \item Shutdown dose rate, coupling a transport-derived spectrum to the
        inventory solve and the emitted photon source back to transport.
  \item Waste classification and recycling limits from the computed
        inventory.
\end{itemize}

\section{Limitations and future work}
\begin{itemize}
  \item One stepper with beginning-of-step rates; predictor-corrector and
        substepping as future work.
  \item Network completeness is bounded by the configured data: a product
        whose parent reaction is absent never appears.
  \item Deep traces are bounds, not significant figures.
  \item No self-shielding or spectrum feedback: the spectrum is an input and
        does not evolve with the inventory.
  \item No uncertainty propagation yet; covariance-driven uncertainties as a
        planned extension.
  \item Charged-particle and photon-induced transmutation channels, and
        further library coverage.
\end{itemize}

\section{Conclusions}
\begin{itemize}
  \item What the code does, what was verified and validated, and at what
        cost in runtime.
  \item Where the design choices (continuous energy, reachability-bounded
        networks, columnar data, one shared workspace with transport) paid
        off.
  \item Availability: package, source, data, documentation and the browser
        deployment.
\end{itemize}

\section*{Back matter}
\begin{itemize}
  \item Code and data availability, with versions and DOIs.
  \item Acknowledgements and funding.
  \item Appendix: CRAM coefficients and the factorization reuse in detail.
  \item Appendix: nuclear data schema tables and library provenance.
  \item Appendix: full V\&V result tables not shown in the main text.
\end{itemize}
